% Standalone manuscript. Compile with pdflatex twice, or tectonic.
% No external bibliography, figures, private documents or shell escape required.
% Prepared for private technical review; not submitted or peer reviewed.
\documentclass[11pt,a4paper]{article}
\usepackage[margin=24mm]{geometry}
\usepackage[T1]{fontenc}
\usepackage[utf8]{inputenc}
\usepackage{lmodern}
\usepackage{amsmath,amssymb,booktabs,array,longtable}
\usepackage{xcolor}
\usepackage{hyperref}
\usepackage{listings}
\setlength{\emergencystretch}{2em}
\hypersetup{colorlinks=true,linkcolor=blue,citecolor=blue,urlcolor=blue}
\lstset{basicstyle=\ttfamily\small,breaklines=true,columns=fullflexible,frame=single}
\newcommand{\sys}{\textsc{Majestic}}
\newcommand{\code}[1]{\texttt{#1}}
\title{Majestic: Evidence-Aware Compilation and Local Deployment of Task-Specific Models}
\author{Majestic Research Project\\
\small Private technical-review manuscript; author and affiliation to be finalized}
\date{23 September 2026}
\begin{document}
\maketitle
\begin{abstract}
Deploying a task-specific model requires more than fitting parameters: the task,
data, resource assumptions, evaluation protocol and executable artifact must agree.
We describe \sys, a research system organizing these decisions around typed
specifications, feasibility predicates, build plans, evidence records and local
artifacts. Its architecture separates specification admissibility, planned
feasibility and artifact evaluation. We report an implementation audit and a
reproducible initial experiment on a deduplicated version of the public UCI SMS
Spam Collection. Across three fixed train/test splits, a TF--IDF nearest-centroid
baseline achieves 94.54\% mean accuracy and 88.94\% macro-F1, while a nearest-neighbor
baseline achieves 97.19\% and 92.99\%, respectively. Packed four-bit centroid
weights reduce weight storage by approximately eightfold but lower mean accuracy
to 92.56\%; storage compression does not ensure inference acceleration. A
separate planted-case experiment agrees with the expected verdicts for 48
declared-effect workflow graphs. Three DistilBERT LoRA runs achieve 98.87\%
mean accuracy and 97.36\% macro-F1. A separate seed-zero QLoRA run with NF4
deployment achieves 99.13\% accuracy in a 73.84 MB deployment directory.
The implementation records provenance separating executed models from
analytical plans. These results establish a reproducible classification and
systems baseline, not universal task coverage, mobile certification, production
safety, or predictive validity of the planner. We specify the experiments needed
to assess those broader hypotheses without treating proposed claims as results.
\end{abstract}

\section{Introduction}
Organizations often need a model to perform a narrow task under practical
constraints: classify a message, extract fields, transform text, retrieve a
grounded answer, or choose an execution path. The desired output is usually an
operational capability with known limitations, rather than a training checkpoint
alone. A successful build therefore has several distinct obligations. The
specification must describe the intended task; the data must support a meaningful
evaluation; the chosen training method must actually execute; deployment must
preserve the evaluated behavior; and any resource commitment must refer to
appropriate evidence.

\sys\ explores a compiler organization for this process. A front end elicits a
typed specification. A planner enumerates supported configurations and records
constraint failures. Data preparation, training and evaluation produce an
artifact together with an evidence manifest. A registry tracks identity and
recall, while a typed directed acyclic graph (DAG) expresses compound workflows.
The useful analogy is an intermediate representation and checked transformations,
not a new universal neural architecture.

This paper makes three bounded contributions. First, it describes the contracts
linking specifications, planned work, executed work and deployment evidence.
Second, it documents concrete implementation failures that can invalidate those
contracts and the corresponding repairs. Third, it supplies reproducible
classification and static-analysis experiments, including unfavorable compression
results. It does not establish priority over all previous systems or claim a new
parameter-efficient learning algorithm.

The current release contains two execution paths. A local studio demonstrates
the specification-to-registry flow using an actual TF--IDF classifier. Its
language-model plan is explicitly analytical. A separate factory executes PEFT
sequence-classification training and HF-format export. Integrating that backend
with every planner recipe remains future work. Treating these paths as a single
fully validated universal factory would overstate the evidence.

\section{Related work}
Compiler infrastructure motivates explicit intermediate representations,
verification and staged lowering. MLIR supports multiple levels of representation
and extensible domain-specific transformations \cite{mlir}. \sys\ applies that
organizational principle to model-building decisions and their evidence; it is
not a replacement for tensor compilers or a demonstrated new compiler theory.

LoRA freezes a pretrained model and learns low-rank parameter updates
\cite{lora}. QLoRA combines low-rank adaptation with a quantized frozen base
\cite{qlora}. These methods are dependencies of the proposed training layer.
The distinction between training-time quantization and deployment quantization
is an execution obligation, not a novel optimization algorithm. The practical
backend follows the PEFT quantization workflow \cite{peft}. DistilBERT is a
compact pretrained transformer suitable for a classification baseline
\cite{distilbert}; adaptation of such a base must not be described as training
a foundation model from scratch.

Hardware-aware specialization is already a substantial research area.
Once-for-All decouples supernetwork training from architecture specialization
under deployment constraints \cite{ofa}. \sys\ instead proposes selecting and
adapting catalogued models for a task and preserving the assumptions behind the
selection. Multi-adapter serving is also prior work: S-LoRA provides specialized
memory management and execution for many concurrent adapters \cite{slora}.
The presence of registry or adapter abstractions in \sys\ does not reproduce
those performance results, especially on mobile devices.

Super-NaturalInstructions demonstrates the breadth of instruction-described NLP
tasks \cite{sni}; it does not validate \sys's proposed primitive taxonomy.
Indirect prompt injection motivates modeling untrusted inputs and privileged
sinks \cite{injection}. Static effect declarations can support conservative
reachability checks, but cannot establish that an arbitrary handler or model is
safe. Finally, binomial uncertainty requires care near boundary probabilities
and at small sample sizes \cite{binomial}. We report Wilson intervals for
classification accuracy rather than treating a point estimate as a guarantee.

\section{System model and contracts}
Let a specification be
\begin{equation}
s=(\tau,\mathcal{I},\mathcal{O},D,r,d,b,q),
\end{equation}
where $\tau$ is a task primitive, $\mathcal{I}$ and $\mathcal{O}$ are input and
output contracts, $D$ identifies data, $r$ records declared rights, $d$ describes
the deployment target, $b$ collects budgets, and $q$ is an acceptance target.
A build plan $p$ records the selected base, adaptation method, hyperparameters,
quantization, runtime and evaluation recipe. An artifact record $a$ binds the
executed model to provenance and evaluation observations.

Three checks have deliberately different meanings:
\begin{align}
G_1(s)&:\text{the request is admissible under the declared policy},\\
G_2(s,p)&:\text{the plan meets modeled constraints},\\
G_3(s,p,a)&:\text{the artifact carries the required evaluation evidence}.
\end{align}
Passing $G_2$ is not a measurement of latency or task quality. Passing a local
quality threshold in $G_3$ is not independent safety certification. The scope
and provenance of each observation must travel with its verdict.

\begin{figure}[ht]
\centering
\small
\fbox{\parbox{0.93\linewidth}{\centering
FORGE $\rightarrow$ Spec IR $\rightarrow$ $G_1$ $\rightarrow$ PLANNER
$\rightarrow$ Build Plan IR $\rightarrow$ $G_2$\\[5pt]
DATA FACTORY $\rightarrow$ TRAINER $\rightarrow$ PROVING GROUND
$\rightarrow$ $G_3$ $\rightarrow$ ARTIFACT\\[5pt]
REGISTRY $\leftrightarrow$ FABRIC local execution and lifecycle controls}}
\caption{Proposed contract structure. The current studio and neural factory
implement different subsets of this flow; their execution provenance is explicit.}
\end{figure}

\subsection{FORGE and the primitive catalogue}
FORGE represents requirements as slots and asks questions when an answer could
change a downstream planning decision, when a required value is missing, or
when the request is contested. Planner sensitivity is useful because uncertainty
in a slot need not imply uncertainty in the chosen plan. However, the utility of
an interview and any attrition coefficient require a human study.

The design documents propose a closed taxonomy, but the implementation and the
newer taxonomy document differ. The implementation includes extraction,
classification, summarization, rewriting, generation, answering, routing and
tool calls. The newer document groups transformation differently and adds
ranking. We therefore treat the taxonomy as provisional. The current local
compiler executes classification and routing-label tasks; unsupported execution
is refused. Neither the catalogue's size nor its presence in code establishes
coverage of real business requests.

\subsection{Planner and resource reasoning}
The planner evaluates candidate configurations against memory, latency,
tokenizer compatibility, data sufficiency, declared licensing restrictions,
offline requirements and build cost. A useful memory accounting identity is
\begin{equation}
M_{\mathrm{total}}=M_{\mathrm{weights}}+M_{\mathrm{KV}}+
M_{\mathrm{runtime}}+M_{\mathrm{adapters}}.
\end{equation}
For a decoder with $L$ layers, $h_{kv}$ key/value heads, head dimension $d_h$,
context length $T$, batch size $B$ and $w$ bytes per cache element,
\begin{equation}
M_{\mathrm{KV}}=2Lh_{kv}d_hTBw.
\end{equation}
These equations require architecture and implementation parameters consistent
with the actual backend. Similarly, document workloads require both prefill
and decode time:
\begin{equation}
t_{\mathrm{request}}\approx n_{\mathrm{in}}/v_{\mathrm{prefill}}
+n_{\mathrm{out}}/v_{\mathrm{decode}}+t_{\mathrm{overhead}}.
\end{equation}
Unmeasured rates are estimates. A profile containing example numbers is marked
as assumed, and cannot become measured merely by passing through an interview.
No mobile latency, thermal or energy result is reported in this study.

The planner's soft quality and cost priors remain unvalidated. A refusal based
on a conservative estimate may reject a build that would succeed. Consequently,
the accuracy of refusal is an empirical question requiring both accepted and
refused configurations to be executed independently of the prediction.

\subsection{Data preparation and training}
The data layer distinguishes real examples, label-preserving augmentation and
synthesis. Additional generated rows are not additional independent human
observations. Splitting and contamination controls must precede augmentation,
and any transformation claimed to preserve a label needs domain-specific checks.

For the public-data factory experiment, normalized duplicate input messages are
removed before a stratified split. A conflicting normalized input is refused by
the general loader. The benchmark ingestion procedure discards entire conflicting
groups and records their number; no such group occurred in the downloaded data.
Vocabulary and inverse-document frequencies are fitted on training inputs only.

The neural backend resolves a Hub reference to a specific revision, uses a
locally loaded tokenizer, fits PEFT adapters and saves the trained adapter.
QLoRA requests use an NF4 base and k-bit preparation. Merging reloads an
unquantized base rather than adding adapter deltas directly to NF4 weights.
Deployment quantization, when selected, is performed separately on the merged
model. Evaluation runs after that operation, and export includes tokenizer and
model files needed for local inference. Unsupported runtime or distillation
requests fail explicitly instead of being recorded as successful no-ops.

\subsection{Evaluation and artifact identity}
Classification accuracy is
\begin{equation}
\hat q=\frac{1}{n}\sum_{i=1}^n\mathbf{1}[\hat y_i=y_i].
\end{equation}
The two-sided 95\% Wilson interval is
\begin{equation}
\frac{\hat q+z^2/(2n)\ \pm\ z\sqrt{\hat q(1-\hat q)/n+z^2/(4n^2)}}
{1+z^2/n},\qquad z=1.96.
\end{equation}
A point estimate above a target with a lower interval endpoint below it is
reported as provisional. This interval concerns the specified sampling model;
it does not remove domain shift, dependence or model-selection bias. Fractional
field-F1 scores must not be rounded into binomial success counts. The local
task axis therefore uses exact-match correctness.

The studio reports seven evaluation axes, including behavioral, regression,
privacy, safety, judge and calibration checks. The present safety/privacy probes
are limited canaries; the default judge is a lexical proxy. They do not justify
claims about rare failure rates or human agreement. Missing model confidence is
reported as unmeasured. Using correctness labels themselves as confidence would
leak evaluation outcomes and manufacture an artificially perfect calibration
score; this behavior was removed during the audit.

For reference predictions $u_i$ and compressed predictions $v_i$, define
\begin{equation}
\phi=\frac{1}{n}\sum_i\mathbf{1}[u_i\ne v_i].
\end{equation}
For exact-match accuracy, $|\hat q(u)-\hat q(v)|\leq\phi$: a row's correctness
can change only when its prediction changes, so summing the per-row absolute
differences proves the inequality. The converse does not hold; improvements and
regressions can cancel. This establishes why flips provide different descriptive
information, not that flips predict field failure better than accuracy delta.

Artifact identity includes actual tensor contents and metadata rather than only
row counts or a planned model name. The cache additionally incorporates supplied
corpus contents, so reusing a data URI cannot silently reuse a different model.
Classifier serving checks the stored tensor digest and lifecycle status. Hashes
detect accidental changes relative to the manifest; without signatures and a
trusted origin they are not remote attestation or protection from an attacker
who can replace both weights and metadata.

\subsection{Fabric and the scope of static guarantees}
Fabric models a workflow as a typed DAG $G=(V,E)$ with declared network effects,
untrusted sources and privileged sinks. A strict policy rejects an active
network node in an offline graph and a tainted path into a privileged operation.
In an adjacency-list implementation, propagation over a topological order needs
$O(|V|+|E|)$ graph operations. The prototype contains repeated edge-list scans,
so that asymptotic bound is a property of the algorithmic formulation, not a
demonstrated implementation performance bound.

Static conclusions are conditional on complete declarations and trustworthy
handlers. The runtime is not a kernel-level network sandbox. A finite output
alphabet may bound channel capacity by $\log_2|\mathcal{O}|$, but even one bit
can authorize a dangerous action. Thus a constrained classifier is not a general
semantic sanitizer. The benchmark uses the strict zero-taint tolerance and does
not test declassification or bounded-capacity authorization policies.

\section{Implementation audit}
The audit identified several ways to obtain plausible reports without the
corresponding evidence. Table~\ref{tab:audit} records the main corrections.
These defects illustrate why interface-level success and passing unit tests are
insufficient for validating a research system.

\begin{table}[ht]
\centering\small
\begin{tabular}{p{0.42\linewidth}p{0.48\linewidth}}
\toprule
Failure mode & Implemented correction \\
\midrule
Assumed Android figures labeled measured & Preserve profile source through elicitation \\
Planned language model reported as executed & Store the actual centroid backend and analytical-plan distinction \\
Correctness used as confidence & Require output-derived confidence or mark unavailable \\
Artifact hash derived from counts & Hash trained tensor contents and validate on load \\
Recalled weights loaded directly & Enforce lifecycle checks during inference \\
Evaluation before quantization & Evaluate and gate the deployed numerical representation \\
QLoRA/compression accepted as no-ops & Execute real backend operations or reject unsupported requests \\
Duplicate inputs crossing splits & Normalize, deduplicate and stratify before fitting \\
Nested runtime packages absent from wheel & Discover and include packages and deployment resources \\
Graph identity omits resource metadata & Bind graph identity to full node declarations \\
\bottomrule
\end{tabular}
\caption{Corrections to executable behavior and reported evidence.}
\label{tab:audit}
\end{table}

\section{Experimental methodology}
\subsection{Data and protocol}
We downloaded the UCI SMS Spam Collection \cite{sms}. The retrieved archive
contains 5,574 parsed rows. Case-folded whitespace-normalized deduplication removes
415 rows, leaving 5,159 unique inputs: 4,517 ham and 642 spam. The archive digest
and exact processed-file digest are recorded in the released local evidence.
This public historical English-language dataset is a test case, not a sample of
the user's production domain.

We use fixed random seeds $0,1,2$ and an 80/20 per-label split, yielding 4,128
training and 1,031 test examples per run. The conditions were fixed before the
benchmark: majority prediction, nearest-centroid and $k=3$ nearest-neighbor
classification. Both learned baselines use training-fitted normalized TF--IDF.
For each learned model, we evaluate float32, signed int8 and packed int4 storage.
This produces 21 reported condition--split results. There is no tuning or repair
against the reported test labels in this experiment.

We report accuracy, macro-F1, confusion matrices, per-split Wilson accuracy
intervals, prediction flip rate and weight-storage ratio. Macro-F1 gives equal
weight to both classes, making the weak majority baseline visible despite its
high accuracy. The three splits overlap, so their average is a descriptive
summary; they are not three independent replications suitable for an unqualified
significance test.

\subsection{Runtime protocol}
The baseline experiment ran locally with Python 3.12.0 and NumPy 2.5.1 on
Windows 11. Training time covers feature fitting and model construction.
Inference timings measure one input per call for the first 32 test messages
after prior batch inference. The table averages per-split medians. It excludes
process startup and file loading, and is neither a service-level percentile nor
a sustained thermal measurement. No phone was used. Quantized NumPy inference
dequantizes stored weights during prediction, an implementation detail that
strongly affects the larger nearest-neighbor matrix.

\subsection{Planted graph experiment}
We construct chains of lengths $1,2,4,8,16,32$ crossed with binary choices for
an untrusted source, a declared network effect and a privileged sink. The 48
cases have expected verdict
\begin{equation}
\mathrm{safe}=\neg\mathrm{network}\ \land\
\neg(\mathrm{untrusted}\land\mathrm{privileged}).
\end{equation}
This is a deterministic planted-case verification experiment. It is not an
independent real-workflow corpus or a comparison against runtime-only detection.

\section{Results}
\begin{table}[ht]
\centering\small
\begin{tabular}{llrrrr}
\toprule
Method & Stored precision & Accuracy (\%) & Macro-F1 (\%) & Flips (\%) & Median ms \\
\midrule
Majority & --- & 87.58 & 46.69 & --- & --- \\
Centroid & float32 & 94.54 & 88.94 & 0.00 & 0.99 \\
Centroid & int8 & 94.44 & 88.80 & 0.23 & 1.08 \\
Centroid & int4 & 92.56 & 85.51 & 2.81 & 1.11 \\
3-NN & float32 & 97.19 & 92.99 & 0.00 & 4.92 \\
3-NN & int8 & 97.12 & 92.82 & 0.06 & 111.81 \\
3-NN & int4 & 97.06 & 92.67 & 0.71 & 217.69 \\
\bottomrule
\end{tabular}
\caption{Means across three fixed splits. Timing is the mean of per-split warm
single-input medians. Flips are relative to the corresponding float32 model.
Weight precision describes storage, not integer-only computation.}
\label{tab:results}
\end{table}

Both learned baselines outperform the majority classifier, particularly in
macro-F1. Four-bit centroid storage loses approximately 1.97 percentage points
of mean accuracy relative to float32 and changes 2.81\% of answers. This is an
observable quality/storage trade-off, not evidence that four-bit precision is
universally optimal. The nearest-neighbor classifier is more accurate on this
dataset but retains a training-example matrix, with correspondingly greater
storage and privacy considerations.

The large latency increases for quantized nearest-neighbor inference expose an
implementation limitation: reconstructing its weight matrix on each call costs
more than the float32 matrix-vector operation. Packed storage therefore cannot
be advertised as an inference-speed improvement in this backend. Caching a
dequantized matrix or using integer kernels would change both timing and resident
memory, and requires a new measurement.

All 48 planted graphs match the expected verdicts. This establishes coverage
of these templates and declared effects only. It does not quantify performance
on real user graphs, prove handler isolation, or establish superiority over
runtime testing.

% NEURAL_RESULTS_BEGIN
\subsection{Pretrained transformer adaptation}
We execute three LoRA runs using the same split seeds $0,1,2$, and one QLoRA
run using seed zero. Each run uses 4,128 training examples and 1,031 test
examples. The base is \code{distilbert/distilbert-base-uncased}, resolved to
the following revision:
\begin{center}\small\ttfamily
12040accade4e8a0f71eabdb258fecc2e7e948be
\end{center}
Training uses three
epochs, batch size 16, learning rate $2\times10^{-4}$, rank 8, scaling parameter
16, adapter dropout 0.05, and query/value projections as the adapter targets.
Inputs are truncated to 128 tokens. The sequence-classification head is trained
alongside the adapters. Hyperparameters are fixed across these runs; neither
early stopping nor model selection uses the reported test labels.

LoRA uses a BF16 base on the GPU. QLoRA uses an NF4 base with double
quantization and gradient checkpointing. Its adapter and classification head
are subsequently loaded onto a fresh unquantized base before merging.
Deployment quantization is a separate NF4 operation, followed by evaluation
and local export. Embeddings and classification-head modules remain outside
the four-bit linear-weight conversion. An initial QLoRA attempt failed while
initializing a quantized classifier head; excluding those head modules fixed
the loader failure before any successful QLoRA training. This was an
implementation correction, not a hyperparameter search on accuracy.

\begin{table}[ht]
\centering\small
\begin{tabular}{lrrrr}
\toprule
Training / deployment & Seed & Accuracy (\%) & Macro-F1 (\%) & Train (s) \\
\midrule
LoRA / unquantized & 0 & 99.13 & 97.99 & 34.03 \\
LoRA / unquantized & 1 & 98.93 & 97.49 & 37.51 \\
LoRA / unquantized & 2 & 98.55 & 96.60 & 37.96 \\
QLoRA / NF4 & 0 & 99.13 & 97.99 & 93.22 \\
\bottomrule
\end{tabular}
\caption{Measured DistilBERT adaptation. Training times are Trainer-reported
loop durations; downloads, loading, merge, quantization and evaluation are excluded.}
\label{tab:neural}
\end{table}

The mean LoRA accuracy is 98.87\% and mean macro-F1 is 97.36\%.
Seed-zero LoRA produces 1,022 correct predictions, with a 95\% Wilson accuracy
interval of $[98.35,99.54]\%$. It is the portable review artifact by a fixed
seed-zero convention, rather than selection of the best observed score.
The unquantized deployment directory occupies 134,867,669 bytes; QLoRA's NF4
deployment occupies 73,841,388 bytes, including tokenizer and configuration.
The latter changes two predictions relative to its own merged unquantized
checkpoint: an answer-flip rate of $2/1031=0.194\%$. These are not flips against
the separate LoRA run. The quantized checkpoint's weight-file compression ratio
is 1.84, not the idealized eightfold ratio of fully quantized float32 weights.

Runs use an NVIDIA GeForce RTX 4060 Laptop GPU, CUDA 12.8, PyTorch
2.11.0+cu128, Transformers 4.57.6, PEFT 0.18.1 and bitsandbytes 0.50.2
under Python 3.12.0 on Windows 11. The saved reports include predictions,
gold labels, training telemetry and actual model revision. Locally initialized
tiny-model integration tests separately check train/save/reload mechanics.
The three split replicates overlap, so their mean is descriptive rather than
three independent datasets. One QLoRA run establishes executable behavior and
a task-specific result, not a general accuracy, memory or speed advantage.
We did not measure neural mobile inference, peak GPU memory or energy.
% NEURAL_RESULTS_END

\section{Research claim ledger}
The architecture motivates a larger research program. Table~\ref{tab:claims}
prevents that program from being mistaken for this paper's results. The thresholds
originate in the supplied plan; this manuscript does not retroactively claim an
external preregistration.

\begin{longtable}{p{0.08\linewidth}p{0.43\linewidth}p{0.40\linewidth}}
\toprule
ID & Proposed hypothesis & Evidence still required \\
\midrule\endhead
C1 & At least 80\% of real requests reduce to the proposed primitives. & Stratified real-request corpus, two independent coders, coverage and inter-rater agreement. \\
C2 & Refusal precision $\geq0.75$ and recall $\geq0.70$. & Execute accepted and refused boundary builds; report per soft predicate without inflating results using trivial hard refusals. \\
C3 & Seed-anchored amplification improves quality by at least five points. & Matched-volume multi-task ablations with fixed data splits and multiple training seeds. \\
C4 & Flip rate adds predictive value for field failure. & Independent field-like labels and model-level analysis beyond accuracy delta; the current flip measurements are descriptive. \\
C5 & Static analysis catches violations runtime testing misses. & Broader adversarial DAG corpus and matched runtime comparator. Current evidence is 48 planted chain cases. \\
C6 & Build-history learning raises first-attempt success by 15 points. & Accumulated real build outcomes, held-out specifications and nearest-neighbor/rule baselines. \\
C7 & Twenty adapters share a mobile base with swaps below 200 ms. & Physical-device adapter, sustained throughput, thermal and battery measurements. \\
\bottomrule
\caption{Unvalidated research hypotheses and their next falsifiable tests.}
\label{tab:claims}
\end{longtable}

\section{Discussion and limitations}
The implementation supports an evidence-oriented workflow, but the evidence
remains narrow. A single historical binary dataset cannot establish extraction,
generation, retrieval, ranking or tool-use quality. Duplicate removal prevents
one form of leakage, while semantically related messages may still cross splits.
Label quality and historical collection biases remain. Repeated inspection of
test outcomes during subsequent development would invalidate their use as an
untouched final evaluation set.

The current studio canary suites are not production safety or privacy audits.
No calibrated model judge or comprehensive regression benchmark was run. Registry
privacy depends on caller identity, storage isolation and operational controls;
local ownership checks alone are insufficient for a multi-tenant service. The
local API is intended for a single user on loopback and is not a hardened public
deployment. Artifact digests provide integrity relative to stored metadata,
not cryptographic proof of trustworthy execution.

Device feasibility combines arithmetic with assumptions. Architecture parameters,
effective precision, concurrency, runtime allocation and workload length can
all change the result. Memory formulas support auditing under their assumptions,
not unrestricted claims of sound refusal. No economic results, user adoption,
willingness to pay, legal determination or competitive advantage were measured.
The business motivation should therefore be evaluated separately from software
correctness and model accuracy.

The full system still requires alignment of the taxonomy, a common execution
contract between planner and neural backend, independent validation of all
supported task recipes, real device measurement, stronger security boundaries
and completion of the research program. These are explicit limitations of the
present release rather than measurements implied by the diagrams.

\section{Reproducibility, data and confidentiality}
The entry point \code{python -m modelrig.review prepare-sms} downloads the public
dataset and records cleaning provenance. \code{python -m modelrig.review benchmark}
recreates the baseline experiment and graph checks. The JSON evidence contains
per-run metrics, sample counts, split digests, tensor digests and environment
information. The standalone LaTeX file contains its own references and figure.

The processed dataset SHA-256 is
\begin{center}\small\ttfamily
80f4c4d8c2ef4d98e51d65911867122643275ece6fd9c33e77f268f43cab21cf
\end{center}
The neural configuration is \code{configs/specs/sms\_lora.yaml}; it fixes the
training seed, split, adapter rank, target modules, learning rate and epoch count.
Model revisions and training telemetry are recorded during execution. Complete
bitwise reproducibility across devices and library versions is not asserted.

The provided confidential patent draft is not a source, dataset or attachment
of this manuscript. Its text was not sent to a research or publication service.
This file is prepared for private review, and no public repository push or
submission is part of the experiment. Author identity and affiliation must be
finalized by the project owner before any submission.

\section{Conclusion}
\sys\ provides an executable setting for investigating whether task-specific
model construction benefits from compiler-style contracts and explicit evidence.
The present results establish reproducible classification baselines and actual
pretrained-transformer adaptation, demonstrate
storage/quality/runtime trade-offs and validate limited declared-effect graph
cases. The audit shows that planned capabilities, measured behavior and reported
certification must be kept distinct. Broader claims about refusal prediction,
task coverage, mobile operation and compounding build-history learning remain
testable hypotheses requiring additional data and experiments.

\begin{thebibliography}{99}
\bibitem{mlir} C. Lattner et al.
\emph{MLIR: A Compiler Infrastructure for the End of Moore's Law}. 2020.
\url{https://arxiv.org/abs/2002.11054}.
\bibitem{lora} E. J. Hu et al.
\emph{LoRA: Low-Rank Adaptation of Large Language Models}. 2021.
\url{https://arxiv.org/abs/2106.09685}.
\bibitem{qlora} T. Dettmers, A. Pagnoni, A. Holtzman, and L. Zettlemoyer.
\emph{QLoRA: Efficient Finetuning of Quantized LLMs}. 2023.
\url{https://arxiv.org/abs/2305.14314}.
\bibitem{peft} Hugging Face.
\emph{PEFT documentation: Quantization}. Accessed September 2026.
\url{https://huggingface.co/docs/peft/developer_guides/quantization}.
\bibitem{distilbert} V. Sanh, L. Debut, J. Chaumond, and T. Wolf.
\emph{DistilBERT, a distilled version of BERT: smaller, faster, cheaper and lighter}. 2019.
\url{https://arxiv.org/abs/1910.01108}.
\bibitem{ofa} H. Cai, C. Gan, T. Wang, Z. Zhang, and S. Han.
\emph{Once-for-All: Train One Network and Specialize it for Efficient Deployment}.
ICLR, 2020. \url{https://arxiv.org/abs/1908.09791}.
\bibitem{slora} Y. Sheng et al.
\emph{S-LoRA: Serving Thousands of Concurrent LoRA Adapters}. 2023.
\url{https://arxiv.org/abs/2311.03285}.
\bibitem{sni} Y. Wang et al.
\emph{Super-NaturalInstructions: Generalization via Declarative Instructions on
1600+ NLP Tasks}. 2022. \url{https://arxiv.org/abs/2204.07705}.
\bibitem{injection} K. Greshake et al.
\emph{Not what you've signed up for: Compromising Real-World LLM-Integrated
Applications with Indirect Prompt Injection}. 2023.
\url{https://arxiv.org/abs/2302.12173}.
\bibitem{binomial} L. D. Brown, T. T. Cai, and A. DasGupta.
Interval estimation for a binomial proportion.
\emph{Statistical Science}, 16(2):101--133, 2001.
\url{https://doi.org/10.1214/ss/1009213286}.
\bibitem{sms} T. A. Almeida and J. M. G. Hidalgo.
\emph{SMS Spam Collection}. UCI Machine Learning Repository, 2011.
\url{https://doi.org/10.24432/C5CC84}.
\end{thebibliography}
\end{document}
